% ==============================================================================
% BAB I. PENDAHULUAN
% ==============================================================================

\chapter{PENDAHULUAN}

\section{Latar Belakang}

Berdasarkan Undang-Undang Nomor 6 Tahun 2014 tentang Desa dan Peraturan Presiden Nomor 39 Tahun 2019 tentang Satu Data Indonesia, pemerintah desa menjadi penyelenggara kegiatan statistik di wilayahnya masing-masing sehingga peran desa sebagai satuan wilayah terkecil menjadi sangat penting. Hal ini karena desa tidak lagi menjadi objek pembangunan, melainkan sebagai subjek dan ujung tombak pembangunan. Oleh karena itu, sebagaimana tertuang dalam Rencana Pembangunan Jangka Menengah Nasional (RPJMN) periode 2020-2026, diperlukan penguatan tata kelola pemerintahan desa dalam upaya pengembangan wilayahnya guna mengurangi kesenjangan dan menjamin pemerataan. Kebijakan desentralisasi dan otonomi daerah menjadi instrumen utama dalam memberikan peluang bagi pemerintah desa untuk membangun desa serta meningkatkan kemandirian dan daya saing desa. Dalam membangun desa, berbagai potensi desa yang dimiliki merupakan modal bagi desa untuk melakukan pembangunan.

Saat ini di desa terdapat berbagai sistem aplikasi (Prodeskel, SDGS Desa, SIK-NG, dst) yang berasal dari berbagai kementerian pusat dan dinas daerah. Sementara, Aparat Desa sebagai narasumber atau produsen data dari berbagai sistem aplikasi tersebut. Dari berbagai sistem yang ada, seharusnya desa memiliki data yang lengkap dan akurat sebagai landasan informasi dalam pengambilan kebijakan pembangunan di desa. Selain itu, permasalahan lainnya adalah mengenai relatif masih rendahnya kualitas dan kapasitas sumber daya manusia (SDM) di pemerintah desa dalam hal pengelolaan data desa. Hal ini berdampak pada rendahnya literasi data di tingkat desa yang pada akhirnya berpengaruh pada komitmen pemerintah desa untuk mengoptimalkan pemanfaatan data dalam kebijakan pembangunan, yang pada gilirannya dapat berdampak pada pengambilan kebijakan yang tidak tepat sasaran.

Data statistik yang dikumpulkan di tingkat desa seharusnya dapat dikelola dan dimanfaatkan oleh pemerintah desa. Selain itu, pengelolaan dan pemanfaatan data desa juga seharusnya selaras dengan prinsip Satu Data Indonesia. Untuk mewujudkannya tidak hanya diperlukan koordinasi dengan penyelenggara kegiatan statistik dan sinkronisasi proses penyelenggaraannya di tingkat desa, tetapi juga diperlukan peningkatan literasi statistik pemerintah desa dalam rangka menjadikan mereka sebagai subjek dalam pengelolaan dan pemanfaatan data di tingkat desa.

Badan Pusat Statistik (BPS) sebagai \textit{leading sector} dalam pengembangan statistik memiliki peran penting dalam peningkatan literasi tersebut. Sebagaimana diamanatkan dalam Undang-Undang Nomor 16 Tahun 1997 tentang Statistik, BPS berkewajiban untuk memberikan pembinaan statistik kepada Kementerian/Lembaga/Satuan Kerja Perangkat Daerah/Institusi Lainnya, termasuk hingga tingkat desa, melalui sistem statistik nasional (SSN) yang berkesinambungan sebagai salah satu bentuk kontribusi dalam peningkatan literasi statistik guna mendukung pembangunan nasional.

Oleh sebab itu, sebagai salah satu perwujudan amanat UU tersebut adalah dengan dilaksanakannya suatu kegiatan pembinaan statistik sektoral di tingkat desa secara berkesinambungan dan komprehensif, yaitu \textbf{Program Pembinaan Statistik Sektoral Desa Cinta Statistik (Desa Cantik)}. Dalam mendukung upaya tersebut, BPS Kabupaten Minahasa Selatan sebagai salah satu Pembina Desa Cantik melaksanakan pembinaan Desa Cantik di Desa Popontolen Kecamatan Tumpaan, Kabupaten Minahasa Selatan.

\section{Gambaran Umum Desa/Kelurahan}

Desa Popontolen termasuk salah satu Desa dalam Kecamatan Tumpaan dengan luas wilayah \textbf{365 ha}. Jumlah penduduk sudah mencapai \textbf{1.570 jiwa} yang merupakan penduduk tetap, dengan jumlah penduduk laki-laki \textbf{816 jiwa} dan perempuan \textbf{754 jiwa}. Luasan wilayah yang begitu potensial saat ini sudah dikembangkan dengan baik oleh masyarakat Desa Popontolen. Terdapat \textbf{569 keluarga}. Keseharian masyarakat Desa Popontolen adalah bertani kesuburan tanah yang luas. Oleh karena itu, hampir sebagian besar penduduk di Desa Popontolen bermata pencaharian sebagai petani karena letak geografis yang sangat mendukung.

\begin{enumerate}[label=\textbf{\arabic*.}]
  \item \textbf{Batas Wilayah/Batas Desa}
  \begin{enumerate}[label=\alph*.]
    \item Sebelah Utara berbatasan dengan Desa Sulu
    \item Sebelah Timur berbatasan dengan Desa Lelema
    \item Sebelah Selatan berbatasan dengan Desa Matani
    \item Sebelah Barat berbatasan dengan Desa Matani
  \end{enumerate}

  \item \textbf{Kondisi Geografis}
  \begin{enumerate}[label=\alph*.]
    \item Ketinggian tanah dari permukaan laut: $\pm$ 3 mdpl
    \item Banyak curah hujan: 179 mm/thn
    \item Topografi: Rendah
    \item Suhu udara rata-rata: 25 -- 30$^\circ$C
  \end{enumerate}

  \item \textbf{Orbitrasi (Jarak dan Pusat Pemerintahan)}
  \begin{enumerate}[label=\alph*.]
    \item Jarak dari pusat pemerintahan kecamatan: 6 km
    \item Jarak dari Ibukota kabupaten Minahasa Selatan: 10 km
    \item Jarak dari Ibukota Provinsi Sulawesi Utara: 43 km
  \end{enumerate}
\end{enumerate}

\section{Gambaran Umum Tata Kelola Desa Cantik di Desa/Kelurahan}

Berdasarkan surat dari BPS RI ketentuan untuk menjadi Desa Cinta Statistik (Desa Cantik) di Tahun 2026 akhirnya BPS Kabupaten Minahasa Selatan melakukan pembentukan Tim Desa Cinta Statistik (Desa Cantik) yang disahkan dengan SK Tim Desa Cinta Statistik (Desa Cantik) Kabupaten Minahasa Selatan. Setelah tim desa cantik terbentuk dilanjutkan dengan kegiatan rapat Tim. Dari hasil rapat dilanjutkan dengan membuat Surat pemberitahuan kegiatan desa cantik ke Pemerintah Daerah Kabupaten Minahasa Selatan untuk menentukan Desa/Kelurahan apa yang akan dijadikan Desa Cinta Statistik (Desa Cantik) di Kabupaten Minahasa Selatan. Dari hasil koordinasi dengan pemerintah daerah akhirnya ditunjuk Desa Popontolen Kecamatan Tumpaan menjadi pilihan dari Pemerintah Daerah. Setelah selesai berkoordinasi dengan Sekda dilanjutkan berkoordinasi dengan Camat Tumpaan dan menyampaikan bahwa Desa Popontolen akan mewakili Desa Cinta Statistik di Kabupaten Minahasa Selatan.

Koordinasi dilanjutkan dengan Hukum Tua Desa Popontolen untuk menyampaikan hasil keputusan mengenai penunjukan Desa Popontolen untuk menjadi Desa Cinta Statistik (Desa Cantik) di tahun 2026 mewakili Kabupaten Minahasa Selatan. Dalam pertemuan juga dibahas mengenai penetapan agen statistik desa cantik sebagai \textit{Brand Ambassador} yang akan lebih banyak berpartisipasi dalam kegiatan desa cantik yang ditunjuk langsung oleh kepala desa (SK terlampir).
